\documentclass[11pt]{article}
\usepackage[margin=0.82in]{geometry}
\usepackage{amsmath,amssymb,amsthm,mathtools,hyperref}
\newtheorem{theorem}{Theorem}
\newtheorem{lemma}{Lemma}
\newcommand{\R}{\mathbb R}
\newcommand{\sos}{\operatorname{SOS}}
\hypersetup{hidelinks,pdftitle={Q222: no fixed polynomial multiplier and the remaining half-line completion}}
\title{No source-independent polynomial Bob multiplier\\for the canonical two-LMI completion}
\author{}\date{}
\begin{document}\maketitle\vspace{-1.8em}
\begin{abstract}
An explicit rational one-parameter family of strictly valid canonical sources
excludes every fixed nonzero nonnegative polynomial Bob multiplier in the
accepted two-LMI completion. The proof uses a polynomial full-state test whose
critical limit is the Motzkin polynomial plus a constant. The required
half-line algebraic cover remains open: its square-root ramification removes
this particular obstruction by an explicit three-square identity. The C1
half-line quadratic form is derived, including its possibly discontinuous
reflection at zero.
\end{abstract}

\section{Statement and the accepted completion identity}
Write the canonical coefficients in the order
$c=(d,\lambda,n,m,C,A,D,\alpha,\beta)$ and put $D_v=1+v^2$.
The four $2\times2$ blocks are
\begin{align*}
P&=\begin{pmatrix}4mv^2&-4mv\\-4mv&4m+(4\beta+2d)D_v\end{pmatrix},&
N&=\begin{pmatrix}4CD_v+4D+2d&(4D+2d)v\\(4D+2d)v&(4D+2d)v^2\end{pmatrix},\\
E&=\begin{pmatrix}4AD_v+(4\alpha+2d)v^2&-(4\alpha+2d)v\\-(4\alpha+2d)v&4\alpha-2dv^2\end{pmatrix},&
K&=4\lambda\begin{pmatrix}1&v\\v&v^2\end{pmatrix}
+4nD_v\begin{pmatrix}0&0\\0&1\end{pmatrix}.
\end{align*}
For a nonnegative polynomial $q\in\R[v]$ and a polynomial $2\times2$ matrix $Z$,
the proposed completion consists of
\begin{equation}\label{eq:blocks}
G_{q,Z}=\begin{pmatrix}q(P+K)&Z\\Z^T&q(N+K)\end{pmatrix}\succeq0,
\qquad
H_{q,Z}=\begin{pmatrix}q(P+E)&q(K-E)-Z\\q(K-E)-Z^T&q(N+E)\end{pmatrix}\succeq0
\end{equation}
for every $v\in\R$.
\begin{theorem}\label{thm:main}
No nonzero $q\in\R[v]$, nonnegative on $\R$ and independent of the source,
admits such a polynomial completion for every valid canonical source.
More precisely, define, for $0<T\leq1/100$,
\begin{equation}\label{eq:source}
\boxed{\begin{aligned}
d&=2T-4T^3,&\lambda=m=A&=T^3,\\
n=D=\alpha&=1-T+T^3,&C=\beta&=(1-T)/2+T^3.
\end{aligned}}
\end{equation}
All nine coefficients are positive, and every source in this family is valid
for all states and all projective measurement angles.
For every fixed $q$ as above there is $\delta_q>0$ such that no polynomial $Z$
satisfies \eqref{eq:blocks} for any source \eqref{eq:source} with
$0<T<\delta_q$.
In particular, the rational sources $T=1/n^2$, $n\geq10$, eventually exclude
any prescribed $q$.
\end{theorem}

We use the established canonical completion identity, without changing its
quantifiers. Let $F_c(u,v;\psi)$ be the cleared canonical slack,
$D_u=1+u^2$, $\psi=(W,X,Y,Z_0)^T$, and
\[
\Phi_c(U,v;a,b)=(1+U)a^TPa+U(1+U)b^TNb
+U(a-b)^TE(a-b)+(a+Ub)^TK(a+Ub).
\]
The two polynomial identities needed here are
\begin{align}
q\Phi_c&=\binom a{Ub}^{\!T}G_{q,Z}\binom a{Ub}
+U\binom ab^{\!T}H_{q,Z}\binom ab,\label{eq:split}\\
u^2D_u F_c(u,v;\psi)&=(L(u)\psi)^T\mathcal P_c(u^2,v)(L(u)\psi),\label{eq:inverse}
\end{align}
where $\mathcal P_c$ is the state matrix of $\Phi_c$ and
\[
L(u)=\begin{pmatrix}u&0&u^2&0\\0&u&0&u^2\\u&0&-1&0\\0&u&0&-1\end{pmatrix}.
\]
These are the accepted source-independent identities. No inverse or cancellation
of a mixed source/testing expression will be used.
If $\deg q=2r$, principal minors of $G_{q,Z}$ imply $\deg Z\leq2r+2$.
Univariate matrix factorization of the two matrices in \eqref{eq:blocks}
then makes \eqref{eq:split}, at $U=u^2$, a polynomial sum of squares whose
roots have degree at most $r+1$ in $v$ and at most two in $u$.
Consequently
\begin{equation}\label{eq:necessary}
q(v)u^2(1+u^2)F_c(u,v;\psi)\quad\hbox{is a polynomial sum of squares.}
\end{equation}
The degree bound depends only on $q$, including when the source or either
completion block loses rank. This consequence of the accepted completion is
all that the obstruction requires.

\section{A strictly valid rational source family}
Let $\eta=1-T$, $w=4-2T+4T^3$ and
\[
f_T=w-\{S+\eta(A_1+B_0)\},\qquad
S=A_0B_0-A_0B_1+A_1B_0+A_1B_1.
\]
In the canonical convention $f=d(1-S/2)+4\sum_i a_i p_i$, its coefficients are
\[
d=4-w,\quad \lambda=m=A=(w-2-2\eta)/4,\quad
n=D=\alpha=(w-2+2\eta)/4,\quad C=\beta=(w-2)/4,
\]
which gives \eqref{eq:source}. This is an exact linear coefficient identity,
verified directly in the accompanying script. Explicitly, in the canonical
qubit basis $A_0=B_0=Z$,
$A_1=((1-u^2)Z+2uX)/D_u$ and $B_1=((1-v^2)Z+2vX)/D_v$;
here $Z=\operatorname{diag}(1,-1)$ and
$X=\left(\begin{smallmatrix}0&1\\1&0\end{smallmatrix}\right)$.
The parent is $D_uD_v$ times the unnormalised state expectation of $f_T$.
Swapping Alice's inputs turns
$f_T$ into the standard symmetric doubly tilted CHSH slack.

Here is a full-state validity proof with explicit constants. Let
$c^2+r^2=e^2+s^2=1$, put $\Omega=1+c+e-ce$, and form the standard Bell matrix
\begin{equation}\label{eq:bell}
H_T=\begin{pmatrix}
\Omega+2(1-T)&(1-c)s&(1-e)r&-rs\\
(1-c)s&-\Omega&-rs&(e-1)r\\
(1-e)r&-rs&-\Omega&(c-1)s\\
-rs&(e-1)r&(c-1)s&\Omega-2(1-T)
\end{pmatrix}.
\end{equation}
Set $p=1-c$, $q_0=1-e$, so $p,q_0\in[0,2]$. Direct expansion, using the
circle equations but neither sine sign, gives
\begin{align}
\det[(4-2T+\sigma)I-H_T]&=\sigma P_T(\sigma;p,q_0)+8pq_0R_T(p,q_0),\label{eq:det}\\
P_T&=\sigma^3+(16-8T)\sigma^2+(84-88T+20T^2)\sigma\notag\\
&\quad+144-240T+112T^2-16T^3+8pq_0(1-T)^2,\notag\\
R_T&=(3+2T-T^2)(p+q_0)-2pq_0-18T+12T^2-2T^3.\notag
\end{align}
Take $\sigma=4T^3$ and $0<T\leq1/100$. Then $P_T\geq144-240T$.
If $R_T\geq0$, the determinant is positive. Otherwise put $a=p+q_0$.
The inequalities $pq_0\leq p+q_0$ and $3+2T-T^2\geq3$ give
$R_T\geq a-18T$, so $a<18T<2$. Since
$2pq_0\leq a^2/2\leq a$, we also have $R_T\geq2a-18T$.
Because $R_T<0$,
\[
8pq_0R_T\geq2a^2R_T\geq4a^3-36Ta^2\geq-432T^3,
\]
where the last step is exactly
\[
4a^3-36Ta^2+432T^3=4(a-6T)^2(a+3T)\geq0.
\]
Thus in this case
\begin{equation}\label{eq:positive-det}
\det(wI-H_T)\geq(144-960T)T^3>0.
\end{equation}

A determinant alone would not prove positivity. At $c=e=1$, the four
eigenvalues of $wI-H_T$ are
\[
4T^3,\quad6-2T+4T^3,\quad6-2T+4T^3,\quad4-4T+4T^3,
\]
all strictly positive. The angle torus is connected and \eqref{eq:det}
never vanishes there. Continuous eigenvalues therefore cannot change sign.
Hence $wI-H_T\succ0$ for every angle pair and both independent sine signs.
This includes literal commuting tests and all projective endpoints.
The accepted real-qubit representation now proves validity of $f_T$ on
all of $Q_{222}$. At $T=0$ the coefficients remain nonnegative and $d=0$;
this endpoint is the immediate positive-event source, also obtained by
continuity. No boundary coefficient has been divided out.

\section{The polynomial full-state test and the limiting contradiction}
Put $T=t^2$, $u=-tx$, $v=ty$ and use the explicit unnormalised state
\begin{equation}\label{eq:state}
\psi_t(x,y)=\begin{pmatrix}
1+\frac23t^4x^2y^2\\
-\frac13t^3x^2y\\
-tx+\frac23t^3xy^2\\
-t^2xy-\frac23t^4x^3y
\end{pmatrix}.
\end{equation}
It is the polynomial numerator of the Alice basis rotation, applied to
$(1,\frac23t^3x^2y,\frac23t^3xy^2,-t^2xy)^T$. Thus no state component
has been omitted, constrained to a limiting contact, or replaced by its
leading value. The canonical source, tests and state satisfy the exact
polynomial expansion
\begin{equation}\label{eq:limit}
F_{c(t^2)}(-tx,ty;\psi_t)
=\frac43t^6\{M(x,y)+2\}+t^8 R(t,x,y),
\qquad
M=x^4y^2+x^2y^4-3x^2y^2+1,
\end{equation}
for a polynomial $R$. The source map, all four state coordinates, and every
coefficient of \eqref{eq:limit} are checked over $\mathbb Q$ by
the included verification script. The file \texttt{CRITICAL\_PULLBACK.txt}
supplies the complete finite expansion.

The leading coefficient also has a short hand calculation. In the standard
frame, use half-angles $tx,ty$ and the state
$\xi=(1,at^3,bt^3,ht^2)^T$. The coefficient of $t^4$ in
$\xi^T(wI-H_{t^2})\xi$ is $4(h+xy)^2$. At $h=-xy$, its coefficient of
$t^6$ is
\[
6\left(a-\frac23x^2y\right)^2+
6\left(b-\frac23xy^2\right)^2+\frac43(M+2).
\]
This determines the displayed state choice while retaining both independent
state-normal squares. The polynomial change to the canonical state and
angle clearing multiply this expectation by $D_u^2D_v$, whose constant
term is one, so they preserve the asserted leading coefficient.

Let $q\ne0$ be nonnegative on $\R$. Its first nonzero Taylor term at the
commuting Bob endpoint is
\[
q(v)=a v^{2k}+O(v^{2k+1}),\qquad a>0.
\]
Suppose completions exist along any sequence $t_j\downarrow0$ in the source
family. Substitute \eqref{eq:state} into \eqref{eq:necessary} and divide by
the positive scalar $t_j^{8+2k}$. Each resulting polynomial in the independent
tests $(x,y)$ is SOS. If $\deg q=2r$, its total degree is at most
$2r+16$, so every SOS root has degree at most $r+8$, independently of $j$.
Its coefficient limit is
\begin{equation}\label{eq:last}
\frac{4a}{3}\,x^2y^{2k}\{M(x,y)+2\}.
\end{equation}
The extra $1+u^2$ factor tends to one. For each fixed real $t_j$, $x,y$
remain arbitrary independent testing variables; the rescaling is used only
to refute the assumed fixed certificate. There is no division by any
expression mixing source and tests.

For completeness, bounded-degree SOS cones are closed for the specific
reason needed here. For a monomial vector $z$ containing all monomials through
the root degree, let
\[
B=\int_{[-1,1]^2}zz^T\,dx\,dy\succ0.
\]
For any Gram matrix $Q_j\succeq0$ representing the $j$th polynomial,
$\operatorname{Tr}(BQ_j)$ is its integral. Coefficient convergence bounds
these integrals, and $B\succ0$ bounds $\operatorname{Tr}Q_j$ and every entry.
A convergent subsequence therefore supplies a PSD Gram matrix for the limit.
This proves that \eqref{eq:last} would be SOS, with no unproved compactness
assumption about a source-dependent function family.

A square monomial can be cancelled from a polynomial SOS. Indeed, if
$x^2h=\sum_i h_i^2$, evaluating at $x=0$ forces each $h_i(0,y)=0$, so each
$h_i$ is divisible by $x$. Repeat for the powers of $y$ and divide the
positive scalar in \eqref{eq:last}. It follows that $M+2$ would be SOS.
This is impossible by an elementary coefficient argument. Since $M+2$ has
total degree six, each possible square root has degree at most three.
On either coordinate axis $M+2=3$, so every root restricts to a constant
there. Each root therefore has the form
$a_0+a_1xy+a_2x^2y+a_3xy^2$. The coefficient of $x^2y^2$ in a sum of their
squares is nonnegative, whereas it is $-3$ in $M+2$.
This proves Theorem~\ref{thm:main}. If there were successful completions
arbitrarily close to $T=0$, they would give the forbidden sequence; hence
all sufficiently small members of the family exclude the chosen $q$.

\paragraph{Stronger separable-multiplier statement.}
The same proof excludes any fixed product $p(u)q(v)$ of two nonzero
polynomials nonnegative on $\R$: for every such pair,
$p(u)q(v)F_{c(T)}$ is not polynomial SOS for all sufficiently small $T>0$.
Indeed, their first nonzero terms at zero are positive even monomials;
under the displayed substitution the leading multiplier is a positive
square monomial in $x,y$, which cancels by the same argument.
The degree of the product is fixed, so any SOS has uniformly bounded root
degree before and after the polynomial substitution. The completion theorem
is the special case $p(u)=u^2(1+u^2)$.

\section{The half-line form and the remaining implication}
The proposed half-line form has an elementary exact justification.
On C1, where $\lambda=n=0$, any polynomial selector satisfying both
unscaled completion inequalities for every $v\geq0$ must have the form
\begin{equation}\label{eq:halfline-selector}
J_+(v)=\begin{pmatrix}
p_1v+p_2v^2&q_2v^2\\
r_0+r_1v+r_2v^2&s_1v+s_2v^2
\end{pmatrix}.
\end{equation}
Indeed, the principal minors of $G_{1,J_+}$ give
$|J_{ij}|^2\leq P_{ii}N_{jj}$, hence degree at most two, even when
the right side vanishes identically. At $v=0$, the zero first diagonal of
$P$ and zero second diagonal of $N$ kill the constant terms of
$J_{11},J_{12},J_{22}$. Finally
$|J_{12}(v)|^2\leq4m(4D+2d)v^4$ kills its linear term.
No positive source weight was assumed or divided out.

The endpoint does not in general force $r_0=0$:
$G_{1,J_+}(0)\succeq0$ gives only
\[
r_0^2\leq(4m+4\beta+2d)(4C+4D+2d),
\]
and, after reversing the sign of $r_0$, $H_{1,J_+}(0)$ adds positive
square blocks to this endpoint quadratic. With
$T_0=\operatorname{diag}(1,-1)$, reflecting by
$J_-(v)=T_0J_+(-v)T_0$ gives, away from zero,
\[
J(v)=\begin{pmatrix}
p_1|v|+p_2v^2&q_2v|v|\\
r_0\operatorname{sgn}(v)+r_1v+r_2v|v|&s_1|v|+s_2v^2
\end{pmatrix}.
\]
Thus two closed half-line charts are legitimate; a global notation must
also retain $\operatorname{sgn}(v)$ unless $r_0=0$ is separately proved.
At zero, averaging the two feasible endpoint selectors supplies the value
with $\operatorname{sgn}(0)=0$. Continuity of the selector is unnecessary.

The obstruction is specific to a polynomial multiplier in the unramified
Bob coordinate. It does not exclude a fixed half-line algebraic cover.
Indeed, the critical polynomial becomes SOS after the single ramification
$y=s^2$:
\begin{equation}\label{eq:halfline}
\boxed{M(x,s^2)=(x^2s^2-1)^2+\bigl(xs^2(s^2-1)\bigr)^2
+2\bigl(xs(s^2-1)\bigr)^2.}
\end{equation}
Adding the constant two preserves this certificate. Equivalently, $\sqrt y$
is one sufficient additional testing function for this limiting obstruction
on $y\geq0$. This supplies a local function, not a unique necessary choice
or a global certificate.

The precise unresolved implication, now for all nine coefficients, is
\begin{equation}\label{eq:open}
\boxed{c\text{ valid}\ \Longrightarrow\
\exists J_+\in\bigl(\R[v]_{\leq2}\bigr)^{2\times2}\ 
\forall v\geq0:\ G_{1,J_+}(v),H_{1,J_+}(v)\succeq0.}
\end{equation}
On C1 this is exactly the eight-parameter question \eqref{eq:halfline-selector}.
If \eqref{eq:open} holds, each quadratic matrix on the half-line has the
fixed-degree factorization
\[
M(v)=(B_0+vB_1)^T(B_0+vB_1)+vC_0^TC_0.
\]
It follows by applying the accepted univariate factorization to $M(s^2)$
and separating even and odd powers. The resulting testing factors are
$1,v,\sqrt v$; rank drops and zero blocks require no inverse.
The missing step is the simultaneous existence of $J_+$ for every valid
source. Its pointwise existence, the degree bound, and the Motzkin identity
do not prove that implication. On the already solved face $d=0$, the
accepted selector $J=K$ remains available, including zero event weights.

No global testing family or affine $Q_{222}$ pencil is obtained here.
The established companion and its pointwise results are not replaced by
this obstruction theorem.
\paragraph{Provenance and reproduction.}
The critical doubly tilted CHSH--Motzkin mechanism is used in
\cite{Chaturvedi}. The present argument proves its own rational upper bound,
canonical source identity, exact polynomial state substitution and
fixed-multiplier obstruction; it uses neither an exact quantum-value formula
nor an identification of NPA certificates with angle-polynomial SOS.
The accepted canonical identities and the unchanged companion are the
project inputs at commit \texttt{918523fc} \cite{Project}.
The accompanying script checks the new identities over $\mathbb Q$ and the
machine-readable source/function data. No SDP computation enters the proof.

\begin{thebibliography}{2}\small
\bibitem{Chaturvedi} A. Chaturvedi, \emph{No Finite NPA Level Characterizes
the Complete Quantum Set in the Simplest Bell Scenario},
\href{https://arxiv.org/abs/2607.14569}{arXiv:2607.14569v1} (2026),
especially Sections S3 and S6--S7.
\bibitem{Project} A. Chaturvedi, \emph{q222-finite-sdp},
\href{https://github.com/anubhav70747/q222-finite-sdp/tree/918523fc022bb7d1b0bf5304fcd485342528456f}{commit 918523fc022bb7d1b0bf5304fcd485342528456f}.
\end{thebibliography}
\end{document}
